\scanpage{209}
\mspage{63 bis}
\uline{\deleted{\unsure{Corollaire}}}\inserted{\unsure{Définition|Définitions}} \unsure{4.6|4.5}. On dit que $u\colon I'\mathbin{\unsure{\leftarrow|-|\to}}I$ dans $\mathrm{Diag}$ est $\mathbb{D}$-\uline{biasphérique}, \deleted{$\mathbb{D}$-\unsure{con}} si elle satisfait les conditions : équivalentes $b,b',b''$ de 4.4, $\mathbb{D}$-\uline{connexe} si elle satisfait la condition $a)$, ou encore (sous réserve d'existence de $u_!$ \deleted{\unsure{et}}\unsure{resp.} de $u_*$) les conditions équivalentes $a'\unsure{)|\}}$, $a''$.

Ainsi, $\mathbb{D}$-biasphérique implique $\mathbb{D}$-connexe, et la réciproque est vraie quand $I$ satisfait \inserted{(p.r. à $\mathbb{D}$)} la condition $3^\circ$ de 4.4. (donc aussi, si elle satisfait l'une des conditions plus faibles $2^\circ,1^\circ$), ou la condition duale de celle-ci (\unsure{car|vu} que les notions $\mathbb{D}$-biasphérique, $\mathbb{D}$-connexe sont autoduales). A cet égard, notons le

\uline{Corollaire 4.7} : Supposons que \inserted{$A=$}$\mathbb{D}(e)$ satisfasse la condition suivante : le foncteur

% slanted
\margin{\uline{NB} Cette condition est satisfaite si $\forall\ \xi,\eta\in A$, on a $\operatorname{Hom}(\xi,\eta)\ne\varnothing$, \unsure{(|[}car on se ramène à l'énoncé : si \deleted{\unsure{les flèches}} $u_1,u_2$ sont deux flèches de $\mathrm{Ens}$ entre ensembles \uline{non vides}, et si $u_1\times u_2$ est iso, $u_1,u_2$ sont iso. ] \deleted{Cela s'applique} Ceci s'applique en particulier si $A$ admet un objet $0$ ; \unsure{ainsi|ou si}, p. ex. si $A$ est additive.}

\[
 \left\{\begin{aligned}
 A\times A&\mathbin{\unsure{\to|-}}A\\
 (\xi,\eta)&\longmapsto \xi\amalg\eta
 \end{aligned}\right.
\]

est conservatif \inserted{(resp. fidèle)}. Alors pour \deleted{\unsure{toute}} flèche \uline{étale} $u\colon J\mathbin{-}I$ dans $\mathrm{Diag}$ \inserted{(et en particulier pour $I_{/i}\mathbin{-}I$)} le foncteur $u_!\colon\mathbb{D}(J)\mathbin{-}\mathbb{D}(I)$ est conservatif (resp fidèle). Si cette condition, ou la condition duale
% archive sheet number: 209
